\chapter*{A note on this critical note}
\addcontentsline{toc}{chapter}{A note on this critical note}

This is a critical note, combining synthesis of the ten books, analysis of the data derived from them, and an evaluation of what the record is and is not. It is a draft of a particular kind. It was made using artificial intelligence, and that has to be said first. The transcription it rests on was made by a language model, Claude Opus 5 and Claude Fable 5.1, reading the Whitney's photographs sheet by sheet in September 2026; the note itself was drafted by Claude Fable 5.1 on 13 September 2026, from those transcriptions and from the data files derived from them, at the request and under the direction of Michel Hua, who set the task and the scope but has not yet read every sentence against the leaves. Eight parallel readings of the transcripts, each by the same model, supplied the quotations and the figures; the prose, the selection and the argument are the model's. A reader should treat the whole as the work of a careful research assistant whose notes have not been checked, which is what it is. It is written from a transcription of the six Hopper ledgers and four of Josephine Hopper's notebooks that was made, sheet by sheet, against the Whitney Museum's photographs, and it is written before any person has checked that transcription against the photographs sheet by sheet. Every quotation below is therefore a quotation of a \emph{reading}, and the reading is dated. Where I quote the leaves I quote them as the transcription gives them, with Josephine Hopper's spelling, her abbreviations and her arithmetic intact, and I give the leaf number she wrote on the paper and the resource reference the museum assigned to the photograph, because the first is what a reader holding the volume can find and the second is the only number that names exactly one image.

The figures are of the same kind. Where I say that the pictures of a year sold for so much, I mean that the rows of the transcribed leaves, read by a reader whose rules are stated, add up to so much; I do not mean that Edward Hopper earned it. The distinction is not modesty. A ledger row with a price in it looks like data in a way a sentence of prose does not, and it will be cited as data whatever caveats stand beside it. So the tables in the appendix carry the date they were generated, and every total in the text is a floor under a number nobody yet knows.

Three readers stand before this one and nothing here claims priority over them. Gail Levin wrote the biography from the diaries and the catalogue raisonné from the record books; Deborah Lyons published the record books in 1997; Madeleine Larson typed the Provincetown diaries. Most of what a first reading of these leaves finds surprising is in Levin already. What this draft can add is not a discovery but a method: the reading placed beside the sheet, so that any sentence in it can be checked against the hand it came out of.

The words of Josephine and Edward Hopper remain in copyright. They are quoted here for the purpose of criticism and study, in the quantities that purpose requires, and this document is a working paper and not a publication.

\vspace{2em}
\noindent\textsc{Conventions.} A citation of the form \lf{Book II}{95}{17982} names the volume, the leaf as numbered on the paper, and the Whitney's ResourceSpace reference for the photograph. Notebooks are cited by page and reference. Quoted matter from the leaves is set between guillemets, \q{like this}, and the writer's own quotation marks inside it are shown as curly double quotes; italics inside a quotation are the writer's own underlining; \ill\ marks a passage the transcriber could not read; a superscript query marks a reading the transcriber offered as doubtful. Struck-through text was struck through in the book. Nothing inside a quotation has been corrected. Every leaf citation is a hyperlink to the transcribed sheet beside its photograph on the edition's site, and every notebook citation to the notebook's reading view. The works discussed are shown as schemas of their composition in main masses only, black and white, at the proportions of the size in Edward Hopper's hand and with a link to the holding museum's own page: the images are in copyright and are not reproduced, and the schemas are the writer's reductions, made from the leaf's description and general knowledge of the picture, not tracings. The figures redraw, from the same data files, the drawings the site makes, and each is signed.
